\paragraph{First-order equivariant class on a compact toy: equivariant lifts and finite-\(N\) truncations (conditional; not closure).}
Replace \(\mathrm{SO}(3,1)\) in the first-order equivariant class above by \(\mathrm{SO}(2)\) and \(\eta\) by the Euclidean contraction, and take the base to be the unit two-sphere, with coframe \(\theta^1=\mathrm d\theta\), \(\theta^2=\sin\theta\,\mathrm d\varphi\), orthonormal frame \(e_1=\partial_\theta\), \(e_2=(\sin\theta)^{-1}\partial_\varphi\) and Levi-Civita connection \(\omega^{12}=-\cos\theta\,\mathrm d\varphi\). On the frame bundle \(Z=\mathrm{SO}(3)\), with fibre angle \(\psi\) and \(\xi=\partial_\psi\), the horizontal lifts \(\hat e_a=e_a-\omega^{12}(e_a)\,\xi\) satisfy exactly \([\hat e_1,\hat e_2]=-\cot\theta\,\hat e_2-\xi\) and \([\xi,\hat e_a]=0\). They do not rotate the frame index, so they violate the first relation of the class for \(\mathcal S_{12}\) any multiple of \(\xi\); read through the second relation alone they would place the anholonomy \([e_1,e_2]=-\cot\theta\,e_2\) of the base frame in the slot of \(\mathsf T\), although the connection is torsion-free. The equivariant lifts \(\mathbb D_a=u_a{}^b\hat e_b\), with \(u_1{}^1=u_2{}^2=\cos\psi\) and \(u_1{}^2=-u_2{}^1=-\sin\psi\), and \(\mathcal S_{12}=\xi\) satisfy both relations exactly, with \(\mathsf T=0\), \(\mathsf F=0\) and \(\mathsf R_{12}{}^{12}=-1\), that is \([\mathbb D_1,\mathbb D_2]=-\mathcal S_{12}\); the opposite orientation of \(u\) gives \(\mathsf T\neq0\), and \(\mathcal S_{12}=-\xi\) violates the first relation. The \(\mathbb D_a\) are divergence-free in \(\sin\theta\,\mathrm d\theta\,\mathrm d\varphi\,\mathrm d\psi\), so \(\mathsf t_a=0\); \(-\mathsf R_{ab}{}^{ae}=\delta_b{}^e\) is the Ricci tensor of the unit sphere; and for every real \(\lambda=\lambda^{12}\) on the base, \([\lambda\mathcal S_{12},\mathbb D_a]=\lambda^c{}_a\mathbb D_c-(\mathbb D_a\lambda)\mathcal S_{12}\), the local rotation of the class. When the curvature is nonzero the commutator of two lifts is not in the span of the lifts over the base functions, so the slot \(\tfrac12\mathsf R_{ab}{}^{cd}\mathcal S_{cd}\) of the class relation is not optional. Its coefficient depends on the basis: for \([\hat e_1,\hat e_2]\) the \(\xi\)-component is \(-1\) in the basis \((\hat e_a,\xi)\) and \(-\csc^2\theta\) in the basis \((e_a,\xi)\), while in the equivariant basis \((\mathbb D_a,\mathcal S_{12})\) it is the constant \(\mathsf R_{12}{}^{12}=-1\). With \(x\) the unit vector of the base point and \(\rho_i=\varepsilon_{ijk}x_j\partial_k\), the lifts \(\tilde\rho_i=\rho_i+c_i\xi\), with \(c_1=-\cos\varphi/\sin\theta\), \(c_2=-\sin\varphi/\sin\theta\) and \(c_3=0\), commute with \(\mathbb D_a\) and \(\mathcal S_{12}\), and \(\sum_i\tilde\rho_i^2=\mathbb D_1^2+\mathbb D_2^2+\mathcal S_{12}^2\). The gauge parameters of \(A_a\) are base functions \(\Lambda\), with \(\mathcal S_{12}\Lambda=0\). On them the ghost operator \([\hat X^a,[\hat X_a,\Lambda]]=-(\mathbb D_1^2+\mathbb D_2^2)\Lambda\) is minus the Laplace--Beltrami operator, with spectrum \(l(l+1)\) and multiplicity \(2l+1\). The redundancy \(\sum_ix_i\rho_i=0\) of the base lifts to \(\sum_ix_i\tilde\rho_i=-\mathcal S_{12}\), which vanishes on base functions.

Two truncations to finite matrices keep different parts of this structure. In the first, the relations are Lie-algebraic here. In the spin-\(j\) representation, \(N=2j+1\), the matrices \(\hat X_a=J_a\) (\(a=1,2\)) and \(\mathcal S_{12}=\mathrm iJ_3\) satisfy both relations with \(\mathsf T=0\), \(\mathsf F=0\) and \(\mathsf R_{12}{}^{12}=-1\) for every \(N\). For the unit four-sphere, \(\hat X_a=\mathrm i\rho(E_{a5})\) and \(\mathcal S_{ab}=\rho(E_{ab})\) satisfy both relations with \(\mathsf T=0\), \(\mathsf F=0\) and \(\mathsf R_{ab}{}^{cd}=-(\delta_a{}^c\delta_b{}^d-\delta_a{}^d\delta_b{}^c)\), so that \(\mathsf R_{12}{}^{12}=-1\), for any finite-dimensional unitary representation \(\rho\) of \(\mathfrak{so}(5)\). The commutative algebra \(\mathcal A\) does not survive at finite \(N\). On a finite-dimensional Hilbert space a commutative \(*\)-algebra is spanned by its minimal projections \(P\). A derivation \(\delta\) that preserves the algebra obeys \(\delta(P)=2P\delta(P)\), hence \(P\delta(P)=0\) and \(\delta(P)=0\). So, for \(\mathcal A\) closed under the adjoint, as the functions on \(Z\) are, every \(\mathrm{ad}\,\hat X_a\) that preserves \(\mathcal A\) vanishes on it; in any case a finite-dimensional commutative algebra has finite spectrum, so no base of positive dimension is realised. Within this, the projection \(P_L\) of \(L^2(Z)=\bigoplus_jV_j\otimes V_j^*\) onto \(j\le L\) commutes with \(\mathbb D_a\), \(\mathcal S_{12}\) and \(\tilde\rho_i\), so the compressions \(P_LM_fP_L\) of the multiplication operators \(M_f\) satisfy \([\hat X_a,P_LM_fP_L]=P_LM_{\mathrm i\mathbb D_af}P_L\). It follows exactly that the relations hold on the range of \(P_L\), that \(\sum_i(P_LM_{x_i}P_L)\,\tilde\rho_i=-\mathcal S_{12}P_L\), and that the ghost operator has the eigenvalue \(l(l+1)\) on the compressed base functions with multiplicity \(2l+1\) for \(l\le2L\). Commutativity fails only at the cutoff: \(P_LM_fP_LM_gP_L-P_LM_{fg}P_L=-P_LM_f(1-P_L)M_gP_L\) vanishes on the blocks \(j\le L-l\) for \(g\) of degree \(l\), and is nonzero on \(j=L\) for \(f,g\) of degree one. The second truncation, the fuzzy sphere, keeps the isometries instead. With \(X_i=J_i/\sqrt{j(j+1)}\) one has \([X_i,X_j]=\mathrm i\hbar_{\mathrm{eff}}\varepsilon_{ijk}X_k\), with \(\hbar_{\mathrm{eff}}=1/\sqrt{j(j+1)}\). The inner derivations \(\mathcal L_i=\hbar_{\mathrm{eff}}^{-1}\operatorname{ad}_{X_i}\) obey the Leibniz rule and \([\mathcal L_i,\mathcal L_j]=\mathrm i\varepsilon_{ijk}\mathcal L_k\) exactly, and \(\sum_i\operatorname{ad}_{X_i}^2/\hbar_{\mathrm{eff}}^2\) on \(\operatorname{End}(\mathbb C^N)\) has the spectrum \(\{l(l+1):0\le l\le N-1\}\) with multiplicity \(2l+1\). By the Casimir identity above, this is the ghost spectrum on base functions, truncated. For odd \(N\) its Hilbert space \(V_j\) is the span of the functions of right weight \(j\) in the block \(j\), which the \(\tilde\rho_i\) preserve and \(\mathbb D_1\) does not, so this truncation carries no frame index. Moving coefficient functions across its product errs from the Poisson bracket at order \(\hbar_{\mathrm{eff}}\) (numerical). The redundancy of the base frame depends on the ordering: on the spin-\(l\) part \(F\) of \(\operatorname{End}(\mathbb C^N)\), \(\sum_iX_i\mathcal L_iF=\tfrac12\hbar_{\mathrm{eff}}\,l(l+1)F\), while \(\tfrac12\sum_i\{X_i,\mathcal L_iF\}=(2\hbar_{\mathrm{eff}})^{-1}[\sum_iX_i^2,F]=0\). Compare the Berezin--Toeplitz finite-size regularization of covariant-derivative configurations, including \(S^2\), of \cite{HattoriSekoTsuchiya2026}. For the Lorentzian fibre the first truncation is not available. \(\mathfrak{so}(3,1)\) is simple and not compact, so every representation of it by anti-Hermitian matrices of finite size is zero, and with \(\mathcal S_{cd}=0\) the first relation at \(a=d\neq c\) gives \(\mathbb D_c=0\). Finite matrices that satisfy both relations with \(\eta\), such as the fundamental representation of \(\mathfrak{so}(4,1)\), have generators that are anti-self-adjoint only for an indefinite form; the same fact for \(\mathfrak{so}(4,1)\supset\mathfrak{so}(3,1)\) underlies the statement in Section~\ref{sec:status-E} that the members with a non-compact five-index algebra exist only for \(N\to\infty\). The toy is compact in base and fibre, Abelian in fibre and two-dimensional. It retains neither \(\mathrm{SO}(3,1)\) nor the \((7,3)\) kinetic signature nor the selection of the class by the matrix equation. It does not derive \(\mathcal A\), the frame or its non-degeneracy, so Assumptions~\ref{ass:frame-lock}, \ref{ass:nondeg} and~\ref{ass:local-lorentz} are not discharged.

\paragraph{The locked frame as a gauge-invariant composite (conditional; not closure).}
Elitzur's theorem \cite{Elitzur1975} rules out a nonzero expectation value of a gauge-variant local field for a compact structure group. Fradkin and Shenker \cite{FradkinShenker1979} showed that, for Higgs fields in the fundamental representation, the Higgs and confinement regimes of such theories are analytically connected without any local order parameter. Both statements are established for compact gauge groups on the lattice; their extension to the noncompact frame group is not established. The expectation value of Equation~\eqref{eq:VEV} is gauge-covariant under Equation~\eqref{eq:Phi-transform}. On a gauge-invariant reading, which Chapters~\ref{ch:vierbein-emergence} and~\ref{ch:lorentz-locking} do not adopt, the locking of Assumption~\ref{ass:frame-lock} would be represented not by the nonzero expectation value itself but by composites that are bilocal in the base, in the spirit of the bilocal fields of \cite{HoKawaiSteinacker2025}, with parallel transporters built from the connection of the class above relating the two points; the composite metric of Equation~\eqref{eq:g-from-e} is invariant under tangent-space rotations. On that reading Assumption~\ref{ass:nondeg} would become a non-degeneracy statement about the two-point composite rather than a consequence of a broken symmetry, and the locking would be read as a complementarity of regimes rather than a local spontaneous breaking. The construction presupposes the connection content that the closure of Problem~D1 in Chapter~\ref{ch:agenda} requires to be specified. It is recorded as a mechanism whose selection by the matrix dynamics is open, and it discharges neither Assumption~\ref{ass:frame-lock} nor Assumption~\ref{ass:nondeg}. Problems~D, F and~N remain open and PROGRAMME-CONDITIONAL.
